\section{Parabolic geometries on tori}

\begin{example}\label{example:TranslationInvariant}
Suppose that $L$ is a Lie group with Lie algebra $\LieL$,
and write the left invariant Maurer--Cartan 1-form
on $L$ as $\ell^{-1}   d\ell$.
Suppose that $G$ is a Lie group and $H \subset G$
is a closed subgroup. Take any linear injection $t : \LieL \to \LieG$
so that the image is complementary to $\LieH$.
Then let $M=L$ and $E=M \times H$. Let $\omega \in \nForms{1}{E} \otimes  \LieG$
be the 1-form
\[
\omega = h^{-1}   dh + \Ad(h)^{-1}\big(t   \ell^{-1}   d\ell\big).
\]
It is easy to check that $\omega$ is the Cartan
connection of Cartan geometry on $L$ modelled on $G/H$.
The group $L$ acts as Cartan geometry automorphisms
by the obvious action. The curvature is
\[
d \omega + \frac{1}{2}\left[\omega,\omega\right]
=
\frac{1}{2}
\Ad(h)^{-1}\left(
\left[
t \ell^{-1}   d\ell,
t \ell^{-1}   d\ell
\right]
-
t
\left[
\ell^{-1}   d\ell,
\ell^{-1}   d\ell
\right]
\right).
\]
In particular, the Cartan geometry is f\/lat if and only if
$t$ is a Lie algebra homomorphism. See~\cite{Hammerl:2007}
for details. Clearly the group $\Aut(L) \times \Aut(G,H)$
acts as isomorphisms of these geometries.
\end{example}


\begin{lemma} Every left invariant Cartan geometry on a
Lie group is equivariantly isomorphic to one
constructed as in Example~{\rm \ref{example:TranslationInvariant}}.
\end{lemma}
\begin{proof}
Take any Cartan geometry $E \to L$ on a Lie group $L$.
Suppose that $L$ acts on $E$ lifting its left action
on itself, commuting with the $H$-action,
and preserving the Cartan connection.
Pick any point $e_0 \in E$. Map
\[
\left(\ell,h\right) \in L \times H \mapsto \ell e_0 h.
\]
Clearly this is an isomorphism of principal bundles,
so from now on we take $E=L \times H$.
Unwinding the def\/inition of a Cartan connection immediately yields
that every Cartan connection on the bundle $L \times H \to L$ has the form
\[
\omega = h^{-1}   dh + \Ad(h)^{-1} \left(\gamma\right),
\]
where $\gamma$ is a 1-form on $L$ valued in $\LieG$,
so that $\gamma+\LieH$ is a linear isomorphism,
i.e.\ $\gamma_{\ell} : T_{\ell} L \to \gtot$ is a
linear injection, for each $\ell \in L$, and
$T_{\ell} L \to \gtot \to \gtot/\gstruc$ is a linear isomorphism.
By translation invariance
\[
\gamma = t \ell^{-1}   d\ell,
\]
for a unique linear map $t$ and
$t \to \gtot \to \gtot/\gstruc$ is a linear isomorphism.
\end{proof}

\begin{theorem}\label{theorem:translation}
Every holomorphic parabolic geometry on any complex torus is translation
inva\-riant, and obtained by the construction of
Example~{\rm \ref{example:TranslationInvariant}}
$($where the Lie group $L$ is the complex torus itself$)$.
More generally, if $M$ is any compact complex
manifold with holomorphically trivial tangent bundle,
then $M=L/\Gamma$ for some discrete
subgroup $\Gamma \subset L$ of a complex
Lie group~$L$. Every holomorphic parabolic geometry
on $M$ is obtained by the construction of
Example~{\rm \ref{example:TranslationInvariant}}
applied to~$L$, and then quotiented by~$\Gamma$.
\end{theorem}

\begin{proof}
If $M$ is a compact complex manifold with holomorphically
trivial tangent bundle, then any basis of the
holomorphic tangent space $T_m M$ at any point $m \in M$
extends to a framing by holomorphic vector f\/ields, say $X_1, X_2, \dots, X_n$.
These then generate a transitive complex Lie group action,
say of a complex Lie group $L$. Brackets of these
holomorphic vector f\/ields can be rewritten in terms
of the vector f\/ields themselves
\[
\left[X_i,X_j\right]=\sum c^k_{ij} X_k,
\]
since the $X_k$ form a basis. The holomorphic functions
$c^k_{ij}$ are constant because $M$ is compact. Therefore $L$ has the same
dimension as $M$, and so $M = L/\Gamma$ \cite{Wang:1954}.

Following the proof of Theorem~\ref{theorem:Reduction}, the structure group of any
holomorphic parabolic geometry $E \to M$ reduces to a reductive group, $G''$
on some subbundle $E'' \subset E$. The Cartan connection~$\omega$ splits into a sum corresponding to the splitting
of $\mathfrak{g}$ into $G''$-invariant subspaces,
and~$\omega''$ (the part valued in $\mathfrak{g}''$)
is a connection form for $E'' \to M$. Take a global
holomorphic coframing on~$M$, i.e.\ a set of linearly independent 1-forms
$\xi^{\alpha}$ forming a basis of each cotangent space of~$M$, for
$\alpha$ varying over noncompact negative roots.
Def\/ine a map $e \in E'' \to h
\in \GL{\gtot/\gstruc}$, by $\omega^{\alpha}=h^{\alpha}_{\beta}
\xi^{\beta}$ (for $\alpha$ and $\beta$ varying over
noncompact negative roots),
and $h(e) =  (h^{\alpha}_{\beta} )$ in the basis~$X_{\alpha}$ for the sum of noncompact negative
root spaces. Under right
$G''$-action,
\[
h\left(r_{g}e\right)=g^{-1}h(e)
\]
for $g \in G''.$ Therefore the quotient map $E \!\to\! \GL{\gtot/\gstruc}/G''$
descends to a map $M \!\to\! \GL{\gtot/\gstruc}/G''$.
The quotient $\GL{\gtot/\gstruc}/G''$ is an af\/f\/ine variety:
see Mumford et al.~\cite[Theorem~1.1, p.~27]{Mumford/Fogarty/Kirwan:1994}
and Procesi~\cite[Theorem~2, p.~556]{Procesi:2007}.
Af\/f\/ine coordinate functions will pull back to functions on $M$,
and therefore must be constant.
Therefore the map $M \to \GL{\gtot/\gstruc}/G''$ is constant.
We have an isomorphism
\[
e \in E'' \to \left(\pi(e),h(e)\right) \in M \times G'',
\]
trivializing the bundle $E''$. We can therefore assume
that $E'' = M \times G''$ and $E = M \times H$.
Again unwinding the def\/inition of a Cartan
geometry,
\[
\omega = h^{-1}   dh + \Ad(h)^{-1}\gamma,
\]
and $\gamma \in \nForms{1}{M} \otimes \LieG$.
The functions $X_i \hook \gamma$
are holomorphic functions on $M$,
so constant, so $\gamma$ pulls back to $L$ to be
$\gamma=t \ell^{-1}   d\ell$ for some constant
linear map $t : \LieL \to \LieG$.
\end{proof}
